\newcommand{\figuraSuperficiesHidrostaticas}{%
\begin{figure}[htbp]
    \centering
    \begin{tikzpicture}[>=Latex,font=\small]
        \node[font=\bfseries] at (6.45,6.05)
            {Superficies de equilibrio hidrostático};

        \draw[very thick,rounded corners=3pt,fill=cyan!4]
            (0.45,0.55) rectangle (6.00,5.45);
        \node[font=\bfseries] at (3.22,5.15) {gravedad uniforme};
        \foreach \y/\tono in {1.25/35,2.05/28,2.85/21,3.65/14,4.45/7}
        {
            \fill[blue!\tono] (0.48,\y-0.38) rectangle (5.97,\y+0.38);
            \draw[blue!60!black,thick] (0.48,\y)--(5.97,\y);
        }
        \draw[ultra thick,green!45!black,-{Latex[length=4mm]}]
            (5.40,4.55)--(5.40,3.45);
        \node[green!40!black,anchor=west] at (5.48,4.02) {$\vec{g}$};
        \node[align=center] at (3.22,0.10)
            {isobaras, isoclinas y equipotenciales\\son planos horizontales};

        \draw[very thick,rounded corners=3pt,fill=orange!3]
            (6.85,0.55) rectangle (12.45,5.45);
        \node[font=\bfseries] at (9.65,5.15) {gravedad central};
        \fill[orange!55] (9.65,2.95) circle (0.28);
        \foreach \r/\tono in {0.85/38,1.40/27,2.00/16}
        {
            \draw[red!70!black,thick,fill=red!\tono,fill opacity=0.20]
                (9.65,2.95) circle (\r);
        }
        \draw[ultra thick,green!45!black,-{Latex[length=3mm]}]
            (11.90,2.95)--(10.75,2.95);
        \node[green!40!black,anchor=south] at (11.35,3.10) {$\vec{g}$};
        \node[align=center] at (9.65,0.10)
            {isobaras, isoclinas y equipotenciales\\son superficies concéntricas};

        \node[align=center,fill=white,rounded corners=2pt,
              draw=gray!35,inner sep=4pt] at (6.45,-0.72)
            {$\vec{\nabla}\rho\times\vec{\nabla}p=\vec{0}$
             \qquad
             $\vec{\nabla}\rho\times\vec{\nabla}\Phi=\vec{0}$};
    \end{tikzpicture}
    \caption{En un fluido barotrópico en equilibrio, las superficies de
    presión, densidad y potencial gravitacional constantes coinciden. Su
    geometría está determinada por el campo gravitacional.}
    \label{fig:superficies-hidrostaticas}
\end{figure}%
}